\documentclass[conference]{IEEEtran}
\IEEEoverridecommandlockouts
\usepackage{booktabs}
\usepackage{multirow}
\usepackage{graphicx}
\usepackage{amsmath}
\usepackage{array}
\usepackage{url}
\usepackage[hidelinks]{hyperref}
\usepackage{xcolor}

\title{Leakage-Aware Evaluation of Subword Supervision Strategies\\
for Cyber Threat Intelligence Named Entity Recognition}

\author{\IEEEauthorblockN{Ammar Aikan, Alaaldeen AL-Shuwaie, and Munif AL-Omisi}
\IEEEauthorblockA{University of Science and Technology--Yemen\\
Sana'a, Republic of Yemen\\
Email address(es)---to be added before submission}}

\begin{document}
\maketitle

\begin{abstract}
Named entity recognition (NER) supports cyber threat intelligence (CTI) analysis by converting unstructured reports into structured mentions of threat actors, tools, campaigns, targets, and related concepts. Reported NER performance, however, can be misleading when duplicated sentences, cross-split overlap, conflicting annotations, malformed records, or invalid BIO transitions remain in a benchmark. This paper audits the DNRTI cybersecurity NER dataset and evaluates how dataset integrity affects conclusions about two BERT subword-supervision strategies. The audit identifies 268 exact token-sequence overlaps between the original training and validation splits, 300 between training and test, and 54 between validation and test. It also finds 23 globally conflicting token-sequence groups, 121 conflict-affected occurrences, malformed input, and invalid BIO sequences. We therefore construct a deterministic, group-aware clean split that quarantines defective groups, preserves all 27 BIO labels and all 13 entity types in each split, and contains zero exact cross-split overlap. Under a controlled single-seed BERT-base-cased experiment, first-subword supervision obtains an exact-span micro-F1 of 0.8675 on the original test split but 0.7002 on the clean test split. BIO-aware subword propagation obtains 0.8622 and 0.7053, respectively. Thus, the apparent 0.0053 advantage of first-subword supervision on the original split reverses to a 0.0051 advantage for BIO-aware propagation on the clean split. Evaluation restricted to originally unseen test sentences substantially narrows the discrepancy, supporting the interpretation that overlap materially influenced the original full-test estimates. The results do not establish universal superiority for either alignment strategy; rather, they demonstrate that dataset integrity can alter both absolute scores and method ranking in CTI NER. We release a reproducible construction policy, cryptographic dataset manifests, exact-span metrics, and explicit limitations to support more defensible evaluation.
\end{abstract}

\begin{IEEEkeywords}
cyber threat intelligence, named entity recognition, BERT, data leakage, dataset quality, subword alignment, DNRTI, reproducibility
\end{IEEEkeywords}

\section{Introduction}
Cyber threat intelligence (CTI) reports contain operationally useful descriptions of adversaries, malware, attack methods, targeted industries, affected regions, and temporal context. Much of this information is published as unstructured prose. Named entity recognition (NER) can reduce the manual effort required to transform such reports into structured records, but a model's measured accuracy is meaningful only when its evaluation data represent genuinely held-out evidence.

Pretrained Transformer encoders such as BERT provide contextual token representations that can be fine-tuned for sequence labeling with limited task-specific architectural change \cite{devlin2019bert}. Their WordPiece tokenization creates an additional design choice for word-level NER corpora: when a labeled word is divided into multiple subtokens, training may supervise only the first subtoken or may propagate a BIO-consistent label over all subtokens. Either choice is plausible. A fair comparison requires identical data, initialization, optimization, and evaluation, and it also requires a clean separation between training and held-out examples.

This separation cannot be assumed. Prior work has shown that train--test overlap can encourage memorization and inflate estimates in natural-language processing \cite{elanggovan2021memorization,lee2022deduplicating}. Recent NER-specific evidence likewise shows that contamination by previously observed entities can overestimate F1 and distort comparisons \cite{babl2025random}. Annotation noise creates a separate but related threat: real labeling errors are not necessarily captured by synthetic-noise experiments and may substantially affect supervised NER \cite{merdjanovska2024noisebench}.

This study examines these risks in DNRTI, a cybersecurity NER dataset containing 13 entity types encoded with 27 BIO labels \cite{dnrti2020}. The work began as a controlled reproduction-oriented comparison within the ACTIT CTI project. Initial evaluation on the supplied splits produced strong full-test results, but a subsequent integrity audit exposed exact duplicates across splits, conflicting annotations for identical token sequences, malformed records, and invalid BIO transitions. These findings motivated a clean, deterministic reconstruction of the dataset and a repeated comparison under the same model family and hyperparameters.

The paper addresses three research questions:

\begin{itemize}
    \item \textbf{RQ1:} What integrity defects are present in the examined DNRTI files, and how extensive are exact cross-split overlaps?
    \item \textbf{RQ2:} How do first-subword and BIO-aware subword supervision compare under matched training conditions?
    \item \textbf{RQ3:} Do conclusions drawn from the original split persist after group-aware cleaning and leakage control?
\end{itemize}

The contributions are fourfold. First, we provide a sentence-level integrity audit that separates parse defects, within-split duplicates, cross-split overlap, conflicting annotations, and invalid BIO sequences. Second, we define a deterministic, model-independent construction policy that quarantines problematic groups and assigns exact token-sequence groups atomically. Third, we conduct a controlled comparison of two subword-supervision strategies using exact-span entity metrics and complementary token metrics. Fourth, we show that cleaning changes both the magnitude of performance and the ordering of the two strategies, while an unseen-only analysis helps distinguish full-test inflation from generalization behavior.

The central claim is deliberately narrow: for the examined DNRTI artifacts and training protocol, conclusions based only on the original full test split are not robust to dataset-integrity controls. We do not claim state-of-the-art performance, statistical superiority, or a causal estimate attributable exclusively to leakage.

\section{Background and Related Work}
\subsection{NER for Cyber Threat Intelligence}
Cybersecurity NER differs from general-domain NER because relevant entities include domain-specific categories such as malicious tools, attacker groups, attack actions, industries, purposes, and sample files. DNRTI was introduced as a large-scale NER dataset for threat intelligence \cite{dnrti2020}. The broader CTI literature combines collection, information extraction, enrichment, and analyst-facing workflows; the ACTIT project that motivated this experiment follows that systems perspective. A recent holistic CTI framework also evaluated BERT and XLNet for CTI entity extraction and reported a BERT F1 of 0.899 \cite{alazab2025holistic}. That value is contextual only here because the published aggregation protocol, corrected dataset version, and full preprocessing details are not sufficiently aligned to support a direct parity claim.

\subsection{BERT and Subword Supervision}
BERT is pretrained through bidirectional contextual objectives and can be adapted to token classification with a task-specific output layer \cite{devlin2019bert}. In word-annotated corpora, WordPiece tokenization may map one source word to several model inputs. A common strategy assigns the word's label to its first subtoken and masks continuation subtokens from the loss. An alternative extends supervision to continuation subtokens. For BIO labels, naive replication is structurally inappropriate when a word begins an entity: the first subtoken retains the B label while later subtokens receive the corresponding I label. We call this second treatment \emph{BIO-aware propagation}.

First-subword supervision avoids overweighting fragmented words, whereas propagation offers denser supervision and explicitly teaches labels for continuation pieces. Their relative behavior can depend on morphology, token fragmentation, label distribution, and evaluation reconstruction. Consequently, the comparison in this paper treats alignment as an empirical factor rather than assuming one universally correct method.

\subsection{Leakage, Duplication, and Annotation Noise}
Random splitting can place identical or highly similar observations on both sides of the train--test boundary. Elangovan et al. distinguish overlap mechanisms across NLP tasks and identify similar input passages and target entities as leakage paths for NER \cite{elanggovan2021memorization}. Lee et al. show more broadly that deduplication reduces memorization and improves the reliability of evaluation \cite{lee2022deduplicating}. Babl et al. provide NER-specific evidence that contaminated entities inflate reported scores and propose splitting designed to reduce contamination \cite{babl2025random}.

Exact sentence overlap is a conservative leakage criterion: it does not capture paraphrases or entity-level contamination, but an exact match is unambiguous and auditable. Annotation conflicts create an additional problem because an identical input may carry multiple target sequences. NoiseBench demonstrates that real NER label noise can be materially harder than simulated noise \cite{merdjanovska2024noisebench}. Our construction therefore avoids silently selecting or rewriting disputed annotations. It quarantines conflicting groups and reports them separately.

\section{Dataset Integrity Audit}
\subsection{Audit Unit and Definitions}
The source files are sentence blocks of token--label pairs. A sentence is \emph{strictly valid} if every nonblank line can be parsed and its label belongs to the expected label map. The identity key is the exact, case-sensitive tuple of parsed tokens. Two occurrences with the same token tuple belong to one group even if their label sequences differ.

We distinguish five conditions:

\begin{enumerate}
    \item \emph{Malformed input}: at least one line in a sentence block cannot be parsed under the frozen format.
    \item \emph{Identical duplicate}: repeated token and label sequences.
    \item \emph{Conflicting duplicate}: identical token sequences associated with distinct label sequences.
    \item \emph{Invalid BIO sequence}: a label transition violates the expected BIO structure.
    \item \emph{Cross-split overlap}: an exact token-sequence group occurs in more than one source split.
\end{enumerate}

The original files were opened read-only; the audit and construction process did not rewrite them. No model metric was used to decide which examples entered a derived split.

\subsection{Original-Split Findings}
Table~\ref{tab:source-audit} summarizes the source files. Raw sentence blocks exceed strictly valid sentences because malformed blocks are rejected as units. Duplicate counts denote occurrences beyond the first within a split.

\begin{table}[t]
\caption{Audit of the Original DNRTI Files}
\label{tab:source-audit}
\centering
\small
\begin{tabular}{lrrr}
\toprule
Measure & Train & Validation & Test \\
\midrule
Raw sentence blocks & 5,251 & 662 & 664 \\
Strictly valid sentences & 5,100 & 655 & 651 \\
Strictly valid tokens & 135,729 & 17,395 & 17,286 \\
Malformed lines & 179 & 7 & 16 \\
Rejected malformed blocks & 151 & 7 & 13 \\
Exact unique token sequences & 3,765 & 612 & 596 \\
Duplicate occurrences & 1,335 & 43 & 55 \\
Conflicting duplicate groups & 17 & 1 & 11 \\
Invalid BIO sentences & 31 & 1 & 4 \\
\bottomrule
\end{tabular}
\end{table}

The within-split counts reveal substantial repetition in training: 1,335 strictly valid occurrences duplicate an earlier exact token sequence. Globally, the audit identifies 23 token-sequence groups with conflicting annotations, affecting 121 occurrences. The conflicts span all 13 entity classes rather than being isolated to one category.

Table~\ref{tab:overlap} reports exact cross-split overlap. The 32 token sequences found in all three source splits are included in the pairwise counts. Pairwise occurrence counts are larger than shared-group counts when a sequence appears repeatedly.

\begin{table}[t]
\caption{Exact Cross-Split Overlap in the Original Files}
\label{tab:overlap}
\centering
\small
\begin{tabular}{lrrr}
\toprule
Split pair & Shared groups & Left occurrences & Right occurrences \\
\midrule
Train--validation & 268 & 493 & 311 \\
Train--test & 300 & 610 & 352 \\
Validation--test & 54 & 66 & 62 \\
\bottomrule
\end{tabular}
\end{table}

These observations invalidate the assumption that the supplied test file is wholly independent of training. They do not prove that every overlapping example is memorized, nor do they quantify approximate or entity-level contamination. They establish a concrete risk that warrants separate full-test and unseen-only reporting.

\section{Leakage-Aware Dataset Construction}
\subsection{Frozen Construction Policy}
The derived dataset uses policy identifier \texttt{dnrti-clean-v2.0}. Construction proceeds as follows:

\begin{enumerate}
    \item Parse source sentence blocks strictly and quarantine malformed blocks.
    \item Group valid observations by their exact, case-sensitive token tuple.
    \item Quarantine groups with multiple distinct label sequences.
    \item Quarantine groups containing invalid BIO transitions.
    \item Retain one canonical instance of each remaining identical group.
    \item Characterize eligible groups by entity-presence strata and sentence-length bins.
    \item Reserve examples so that every entity type occurs in train, validation, and test.
    \item Allocate groups deterministically with seed 42, never dividing one group across splits.
\end{enumerate}

The policy does not normalize BIO labels and does not adjudicate conflicts. This conservative choice sacrifices data volume but prevents undocumented corrections from becoming hidden experimental degrees of freedom.

\subsection{Derived Dataset}
The process yields 4,334 eligible unique groups. Table~\ref{tab:clean-splits} describes the derived files. Every clean test sequence is absent from clean training by construction. SHA-256 values bind the reported results to exact file contents.

\begin{table*}[t]
\caption{Leakage-Aware DNRTI Clean-v2 Splits}
\label{tab:clean-splits}
\centering
\small
\begin{tabular}{lrrl}
\toprule
Split & Sentences & Tokens & SHA-256 \\
\midrule
Train & 3,460 & 89,343 & \texttt{50fd29499e4e1e324b6b43bd765d48064836ab831ef2d55262968ee10aba93a6} \\
Validation & 436 & 13,040 & \texttt{90440d271e6476f019fe6c90ca288b07ec7b2bbfa8d5972ec316a6c692907325} \\
Test & 438 & 11,205 & \texttt{f5d80c3ef68b8ff6f2b49f28fee28ab3a6d60f0f0e99c6ef64f2ab5d2e425359} \\
\bottomrule
\end{tabular}
\end{table*}

Eight pretraining integrity gates pass: zero train--validation, train--test, and validation--test exact overlap; zero conflicting duplicate annotations; zero malformed lines; zero invalid BIO transitions; preservation of the exact 27-label map; and representation of every entity type in each split.

\section{Experimental Method}
\subsection{Model Variants}
Both variants initialize from the same fresh \texttt{bert-base-cased} checkpoint and add a token-classification head for the 27-label map.

\textbf{First-subword supervision (FS).} The first WordPiece associated with a source word receives the word's BIO label. Continuation pieces receive the ignore index $-100$ and do not contribute to the training loss. Predictions are reconstructed at word level from first pieces.

\textbf{BIO-aware propagation (BA).} The first piece receives the source label. Continuations of an entity word receive the matching I-type label; continuations of O words remain O. This supplies loss targets to all lexical pieces while respecting entity continuation semantics. Word-level evaluation uses the reconstructed source-token sequence.

In the initial original-split experiment, propagation added 44,419 supervised training positions relative to first-subword masking. The original training split contained 19,517 fragmented words (13.906\% of parsed words) and 44,419 continuation pieces. No sequence was truncated at the configured maximum length in that run.

\subsection{Controlled Configuration}
Table~\ref{tab:hyperparameters} lists the shared configuration. The effective training batch is $2\times8=16$. Both models train on CPU with four PyTorch threads. Validation exact-span micro-F1 selects the best epoch. Apart from alignment, all listed factors are held constant.

\begin{table}[t]
\caption{Controlled Training Configuration}
\label{tab:hyperparameters}
\centering
\small
\begin{tabular}{ll}
\toprule
Parameter & Value \\
\midrule
Base model & BERT-base-cased \\
Epochs & 4 \\
Physical batch size & 2 \\
Gradient accumulation & 8 \\
Effective batch size & 16 \\
Evaluation batch size & 4 \\
Learning rate & $1\times10^{-4}$ \\
Optimizer & AdamW \\
Weight decay & 0.01 \\
Adam epsilon & $1\times10^{-12}$ \\
Maximum gradient norm & 1.0 \\
Maximum sequence length & 120 \\
Random seed & 42 \\
PyTorch CPU threads & 4 \\
\bottomrule
\end{tabular}
\end{table}

\subsection{Metrics}
The primary outcome is strict entity-level micro-F1. A predicted entity is a true positive only when its type and complete word-level span match a gold entity exactly. Let $TP$, $FP$, and $FN$ denote matched predictions, unmatched predictions, and missed gold entities. Then
\begin{equation}
P=\frac{TP}{TP+FP},\quad R=\frac{TP}{TP+FN},\quad F_1=\frac{2PR}{P+R}.
\end{equation}
We additionally report macro-F1 across the 13 supported entity types, token accuracy, and micro token-F1 excluding O. Token accuracy is secondary because O dominates the sequence. The original split is also evaluated on the subset of test sentences whose exact token sequences are absent from original training (the \emph{unique-unseen} subset).

\subsection{Research Integrity Controls}
Construction is completed before training and is independent of model results. All derived files and trained artifacts have recorded SHA-256 hashes. The source dataset remains unchanged. The experiment does not access production CTI records, credentials, or private operational data. Neither experimental model was promoted to ACTIT production; operational deployment decisions are outside the scientific comparison.

\section{Results}
\subsection{Original Versus Clean Evaluation}
Table~\ref{tab:main-results} provides the central comparison. On the original full test split, FS has a 0.0053 F1 advantage over BA. On Clean-v2, BA has a 0.0051 advantage. The ranking therefore reverses, although both differences are too small to support a significance claim from one seed.

\begin{table*}[t]
\caption{Exact-Span Entity Results Across Evaluation Regimes}
\label{tab:main-results}
\centering
\small
\begin{tabular}{llrrrrrr}
\toprule
Dataset/evaluation & Strategy & Precision & Recall & Micro-F1 & Macro-F1 & Token acc. & O-excl. token F1 \\
\midrule
Original full test & First-subword & 0.8527 & 0.8829 & 0.8675 & 0.8814 & 0.9630 & 0.8906 \\
Original full test & BIO-aware & 0.8480 & 0.8769 & 0.8622 & 0.8739 & 0.9606 & 0.8846 \\
Original unseen-only & First-subword & 0.6966 & 0.7441 & 0.7196 & -- & -- & -- \\
Original unseen-only & BIO-aware & 0.7074 & 0.7492 & 0.7277 & -- & -- & -- \\
Clean-v2 test & First-subword & 0.6717 & 0.7312 & 0.7002 & 0.6353 & 0.9116 & 0.7387 \\
Clean-v2 test & BIO-aware & 0.6799 & 0.7326 & 0.7053 & 0.6407 & 0.9117 & 0.7420 \\
\bottomrule
\end{tabular}
\end{table*}

Relative to original full-test evaluation, clean-test F1 is lower by 0.1673 for FS and 0.1570 for BA. The original full-to-unseen gaps are 0.1480 and 0.1345, respectively. In contrast, the differences between original unseen-only and Clean-v2 are 0.0194 for FS and 0.0225 for BA. This pattern is consistent with overlap inflating the original full-test estimates. It is not a pure causal estimate because Clean-v2 also removes malformed, conflicting, and invalid groups, deduplicates repeated inputs, reallocates the eligible groups, and reduces training size.

On Clean-v2, BA predicts 1,559 entities and matches 1,060, producing 499 false positives and 387 false negatives. FS predicts 1,575 and matches 1,058, producing 517 false positives and 389 false negatives. BA's small aggregate advantage therefore comes from two additional matches and 18 fewer false positives.

\subsection{Per-Type Behavior on Clean-v2}
Table~\ref{tab:per-type} reports exact-span F1 by entity type. Performance varies substantially by class. Both strategies perform strongly on \textit{SecTeam}, \textit{Time}, and \textit{Area}, but weakly on \textit{Purp}. BA improves seven classes (Features, Idus, OffAct, Org, SamFile, Tool, and Way) and declines on six. The largest positive BA difference is for Way (+0.0409), whereas its largest negative difference is for Exp ($-0.0434$). These heterogeneous changes caution against interpreting the small micro-average difference as uniform behavior.

\begin{table}[t]
\caption{Clean-v2 Exact-Span F1 by Entity Type}
\label{tab:per-type}
\centering
\small
\begin{tabular}{lrrr}
\toprule
Type & Support & FS F1 & BA F1 \\
\midrule
Area & 159 & 0.8199 & 0.7951 \\
Exp & 28 & 0.8070 & 0.7636 \\
Features & 34 & 0.4242 & 0.4918 \\
HackOrg & 321 & 0.7486 & 0.7278 \\
Idus & 81 & 0.6494 & 0.6667 \\
OffAct & 158 & 0.7152 & 0.7390 \\
Org & 93 & 0.5426 & 0.5775 \\
Purp & 26 & 0.2821 & 0.2424 \\
SamFile & 41 & 0.4533 & 0.5000 \\
SecTeam & 104 & 0.8835 & 0.8358 \\
Time & 135 & 0.8540 & 0.8370 \\
Tool & 220 & 0.6426 & 0.6753 \\
Way & 47 & 0.4364 & 0.4773 \\
\bottomrule
\end{tabular}
\end{table}

\subsection{Validation Selection}
FS reaches its best validation entity F1, 0.6869, at epoch 4. BA reaches 0.6868 at epoch 3. Clean-test F1 differs by approximately half a percentage point despite nearly identical best validation F1 values. This reinforces the need to avoid post-hoc claims from small differences and a single random seed.

\section{Discussion}
\subsection{Answering the Research Questions}
\textbf{RQ1.} The examined DNRTI files contain several distinct integrity risks. Exact overlap is extensive enough that 300 unique token sequences occur in both original training and test. Repetition within the training file is also material, and some identical inputs carry conflicting target labels. Malformed records and invalid BIO transitions further complicate implicit preprocessing. A headline score without a precise parser and split audit is therefore difficult to interpret or reproduce.

\textbf{RQ2.} Under matched conditions, neither subword strategy dominates consistently. FS is slightly better on the original full test, whereas BA is slightly better on original unseen-only evaluation and Clean-v2. Per-class results are mixed. The defensible conclusion is not that propagation is universally superior, but that strategy rankings of this magnitude are sensitive to evaluation composition.

\textbf{RQ3.} The original conclusion does not persist after integrity controls. Full-test F1 decreases substantially, and the direction of the alignment comparison reverses. The clean scores are much closer to the original unseen-only scores than to original full-test scores. This provides converging, but not causal, evidence that overlap and dataset defects influenced the optimistic full-test estimates.

\subsection{Why Token Accuracy Is Insufficient}
Both clean models exceed 0.911 token accuracy, yet their exact-span micro-F1 is near 0.70 and macro-F1 near 0.64. O-excluded token-F1 is also lower than overall accuracy. In CTI extraction, a partially identified multi-token entity can be operationally wrong even when most O tokens are classified correctly. Exact-span metrics should therefore be the primary evidence, accompanied by class-specific analysis.

\subsection{Implications for CTI NER Evaluation}
The study suggests a practical evaluation sequence for CTI NER:

\begin{enumerate}
    \item audit parsing and BIO validity before training;
    \item group exact duplicates before split assignment;
    \item isolate conflicting annotations rather than silently selecting one;
    \item report exact cross-split overlap and unseen-only performance;
    \item bind datasets and model artifacts to hashes;
    \item report exact-span, O-excluded, and per-class metrics;
    \item treat small single-seed differences as descriptive, not conclusive.
\end{enumerate}

These controls are inexpensive compared with model training and can prevent an apparent architectural gain from being mistaken for improved generalization.

\subsection{Relationship to ACTIT}
The experiment was conducted as part of ACTIT, a broader CTI platform. Its role is methodological: to determine whether candidate DNRTI-BERT models have evidence strong enough for integration. Scientific training success and production promotion are separate decisions. The clean models did not satisfy the project's predefined entity-F1 gates, and production was not changed. This separation avoids presenting an offline result as deployed behavior and illustrates how leakage-aware evaluation can serve as an operational safety gate.

\section{Threats to Validity and Limitations}
\textbf{Single seed and no uncertainty intervals.} Each condition uses seed 42. The small differences between strategies may be explained by training variance. Multiple seeds and paired resampling are necessary before testing statistical superiority.

\textbf{One encoder and two alignment strategies.} The experiment covers BERT-base-cased only. Results may not transfer to domain-adapted encoders, other tokenizers, CRF decoders, larger models, or generative extraction approaches.

\textbf{Cleaning changes several factors.} Clean-v2 simultaneously removes malformed blocks, conflicting groups, invalid BIO sequences, and exact duplicates, and it reallocates groups. Its training set is smaller and its distribution differs from the original. The observed decrease cannot be attributed exclusively to train--test overlap.

\textbf{Exact overlap is a lower bound.} The audit detects exact, case-sensitive token-sequence identity. It does not detect paraphrases, near duplicates, shared document context, or entity-level contamination. Residual semantic leakage may remain.

\textbf{No manual conflict adjudication.} Quarantine avoids arbitrary rewriting but may discard informative examples. Expert relabeling could recover data and produce a stronger gold standard.

\textbf{Reference-score non-equivalence.} A previously reported BERT score of 0.899 \cite{alazab2025holistic} is not treated as a directly comparable baseline because the precise corrected data, aggregation, preprocessing, and evaluation implementation are unavailable to this study.

\textbf{Artifact availability.} The Clean-v2 construction/training path and result artifacts are present in the ACTIT repository. An earlier original-split training utility was not retained in the final worktree, although its metrics, checkpoints, provenance, and manifests remain. This limits exact end-to-end reproduction of that earlier run and is disclosed rather than obscured.

\textbf{Dataset distribution and licensing.} The hashes and policy support verification, but redistribution of derived data depends on the original dataset's license. A public release must verify those terms or publish reconstruction code without repackaging restricted content.

\section{Reproducibility Statement}
The Clean-v2 experiment records the source and derived file hashes, construction-policy version, seed, model checkpoint identity, command-line configuration, training history, exact-span outputs, token metrics, per-class metrics, artifact hashes, and final deployment decision. The training entry point is \texttt{ml/training/bert\_dnrti\_clean\_v2.py}; the experiment root is \texttt{ml/experiments\_clean\_v2/20261014T110000Z}. The model artifact hashes are \texttt{d6b2b081...890b37b} for FS and \texttt{734c36e0...ca7ff5} for BA. The execution environment used Python 3.12.3, CPU-only PyTorch 2.14, Transformers 5.17, Tokenizers 0.23.2, NumPy 2.5.1, and scikit-learn 1.9 on an AMD EPYC CPU. Paths and versions are reported as provenance, not as requirements that guarantee numerically identical execution on every platform.

To support independent verification, a release should include: (1) the construction utility and frozen policy; (2) source-file hashes rather than source contents when licensing is uncertain; (3) a generated manifest with group assignments; (4) the two training commands; (5) metric implementation tests; and (6) machine-readable results. The authors intend to publish the paper, evidence map, and legally redistributable reproducibility artifacts through the ACTIT project website in addition to the archive or supplementary-material mechanism required by the selected venue. The permanent public URL and release version will be inserted only after publication infrastructure and dataset-license compliance have been verified.

\section{Conclusion}
This paper audited DNRTI and evaluated two BERT subword-supervision strategies under original and leakage-aware conditions. The supplied splits contain exact cross-split overlap, within-split repetition, conflicting annotations, malformed records, and invalid BIO transitions. A deterministic group-aware reconstruction removes these observed defects while preserving the complete label map and entity-type coverage.

The empirical conclusion changes after cleaning. First-subword supervision leads by 0.0053 F1 on the original full test, but BIO-aware propagation leads by 0.0051 on Clean-v2. Both clean results are far below the original full-test scores and much closer to results obtained on originally unseen sentences. Dataset integrity therefore affects not only reported performance but also the relative ordering of modeling choices.

Future work should repeat training across multiple seeds, use paired confidence intervals, separate individual cleaning ablations, inspect entity-level and near-duplicate contamination, manually adjudicate quarantined conflicts, and evaluate domain-adapted encoders on an external CTI corpus. Until such evidence is available, leakage-aware dataset construction and transparent limitations are more defensible than claims of model superiority.

\section*{Acknowledgment}
The authors thank Dr. Raed Al-Dhabhani for supervising the ACTIT graduation project at the University of Science and Technology--Yemen. Any additional institutional, infrastructure, funding, or dataset acknowledgment will be added only after confirmation.

\bibliographystyle{IEEEtran}
\bibliography{references}

\end{document}
